\chapter{Reproducibility and Packaging}
\label{ch:reproducibility}

\newthought{A scan you cannot reproduce} is hard to trust and hard to share. \Jarvis\ is built so a
run is self-contained by default, and so you can bundle it for a collaborator or an archive. This
closing chapter of the part pulls those threads together.

\section{What makes a run reproducible}
\label{sec:repro-what}

Several design choices, met throughout the book, add up to reproducibility:

\begin{description}[style=nextline, leftmargin=1.2em]
  \item[Project-local everything] the task card, dependencies, and outputs all live under the
        project root, addressed with \marker{\&J/} (Chapter~\ref{ch:core-concepts}); nothing is
        scattered across the machine.
  \item[Declared environment] \yamlkey{EnvReqs} records the platform and packages a run needs and
        checks them before starting (Chapter~\ref{ch:envreqs}).
  \item[Captured configuration] checkpoints store the exact config the run used, so a resumed run
        stays consistent with itself (Chapter~\ref{ch:checkpoint}).
  \item[Replayable points] the \sampler{CSV} sampler re-runs an exact list of points
        (Chapter~\ref{ch:basic-samplers}), and many samplers store a random seed.
\end{description}

\section{Packaging a project}
\label{sec:repro-pack}

\cmd{Jarvis project pack} (Chapter~\ref{ch:project-cli}) bundles a project, with the mode flag
choosing how much travels (Table~\ref{tab:pack-modes}).

\begin{table}[ht]
  \footnotesize
  \begin{tabular}{@{}llp{0.46\linewidth}@{}}
    \toprule
    \textbf{Mode} & \textbf{Flag} & \textbf{Use it to\ldots} \\
    \midrule
    Share    & \cli{-{}-share} (default) & send results to a collaborator (lighter) \\
    Repro    & \cli{-{}-repro}           & let someone unpack and re-run the scan \\
    Full     & \cli{-{}-full}            & archive everything for the long term \\
    \bottomrule
  \end{tabular}
  \caption{Packaging modes.}
  \label{tab:pack-modes}
\end{table}

\begin{shellbox}
Jarvis project pack . --repro
\end{shellbox}

\noindent The \cli{-{}-repro} package is the one to choose when reproducibility is the goal: it
carries what is needed to unpack and run the scan again, not just the finished numbers.

\begin{jtip}
Write a manifest first with \cli{-{}-man}, then edit its \yamlkey{include}/\yamlkey{exclude} lists
to keep large build trees or raw outputs out of the bundle, and pack from that manifest. You get a
reproducible package without shipping gigabytes of intermediate files.
\end{jtip}

\section{Good habits}
\label{sec:repro-habits}

A few practices make future-you (and your collaborators) grateful:

\begin{itemize}
  \item Keep the package payloads under \file{deps/} so the project carries its own inputs.
  \item Pin versions in \yamlkey{EnvReqs} so a mismatch is caught early, with a clear message.
  \item Record a seed where the sampler supports one, and keep the points of interest as a
        \textsc{csv} you can replay.
  \item Note which \Jarvis\ release produced a result---\cmd{Jarvis -{}-version} prints it, and it
        is captured in the run summary and logs.
\end{itemize}

\begin{jnote}
\Jarvis\ evolves between releases. Recording the version that produced a result---and packaging the
project with \cli{-{}-repro}---is the surest way to reproduce it later, even as the framework moves
on.
\end{jnote}

\section{Summary}
\label{sec:repro-summary}

Reproducibility in \Jarvis\ is mostly free: project-local layout, a declared environment, captured
config, and replayable points. When you need to share or archive, \cmd{Jarvis project pack}
---especially \cli{-{}-repro}---turns a project into a portable, re-runnable bundle.
